% 共享导言区。与正文分开，使正文保持可读。
\usepackage[UTF8]{ctex}
\setCJKmainfont{Noto Serif CJK SC}
\setCJKsansfont{Noto Sans CJK SC}
\setCJKmonofont{Noto Sans Mono CJK SC}
\usepackage{amsmath,amssymb,amsthm,mathtools}
\usepackage{booktabs,multirow,array}
\usepackage{algorithm}
\usepackage[noend]{algpseudocode}
\usepackage{graphicx}
\graphicspath{{../../results/leak/}}
\usepackage[table]{xcolor}
\usepackage{geometry}
\geometry{a4paper,margin=2.4cm}
\usepackage[colorlinks=true,linkcolor=blue!55!black,citecolor=blue!55!black,
            urlcolor=blue!55!black]{hyperref}
\usepackage{caption}
\captionsetup{font=small,labelfont=bf}
\usepackage{enumitem}
\setlist{itemsep=2pt,parsep=0pt,topsep=3pt}
\usepackage{fancyhdr}
\setlength{\headheight}{14pt}
\addtolength{\topmargin}{-2pt}
\pagestyle{fancy}\fancyhf{}
\fancyhead[L]{\small 算子谱先验}
\fancyhead[R]{\small\thepage}
\renewcommand{\headrulewidth}{0.4pt}

% 中文算法环境
\floatname{algorithm}{算法}
\renewcommand{\algorithmicrequire}{\textbf{输入：}}
\renewcommand{\algorithmicensure}{\textbf{输出：}}
\algrenewcommand\algorithmicfor{\textbf{对}}
\algrenewcommand\algorithmicdo{\textbf{执行}}
\algrenewcommand\algorithmicwhile{\textbf{当}}
\algrenewcommand\algorithmicif{\textbf{若}}
\algrenewcommand\algorithmicthen{\textbf{则}}
\algrenewcommand\algorithmicelse{\textbf{否则}}
\algrenewcommand\algorithmicend{\textbf{结束}}
\algrenewcommand\algorithmicreturn{\textbf{返回}}

% 记号
\newcommand{\Lop}{\mathcal{L}}
\newcommand{\Aop}{\mathcal{A}}
\newcommand{\Sop}{S_{\mathrm{op}}}
\newcommand{\bxi}{\boldsymbol{\xi}}
\newcommand{\NRMSE}{\mathrm{NRMSE}}
\newcommand{\Om}{\Omega}
\newcommand{\norm}[1]{\lVert #1 \rVert}
\newcommand{\code}[1]{\texttt{\small #1}}

\theoremstyle{definition}
\newtheorem{principle}{准则}
\theoremstyle{plain}
\newtheorem{proposition}{命题}

% 结论框：把"能守住的话"和正文分开
\usepackage{tcolorbox}
\tcbuselibrary{skins,breakable}
\newtcolorbox{takeaway}[1][]{
  enhanced, breakable, colback=blue!3, colframe=blue!45!black,
  boxrule=0.5pt, left=6pt, right=6pt, top=5pt, bottom=5pt,
  fonttitle=\bfseries\small, title={#1}}
